\homework{2}{Sun 29 Sep 2024 01:07}{}

\begin{enumerate}
\item Let $Q_k$ be a Whitney Cube center at $y^k$ and set $Q_k^*=2 \sqrt{n} Q_k$. Let $\varepsilon>0$ be given. Prove that
1. If $x \notin Q_k^*$, then $\left|x-y^k\right| \geq 2\left|y-y^k\right|$ for any $y \in Q_k$.
2. If $x \notin Q_k^*$ and for some $y_0 \in Q_k,\left|x-y_0\right|=\varepsilon$, then there is a constant $C_1$ depending only on $n$ such that $Q_k \subset \overline{B\left(x, C_1 \varepsilon\right)}$.
3. If $x \notin Q_k^*$ and for some $y_0 \in Q_k,\left|x-y_0\right|=\varepsilon$, then there is a constant $C_2$ depending only on $n$ such that $|x-y| \geq C_2 \varepsilon$ for all $y \in Q_k$.
\begin{proof}
(1)
Since $x \in Q_k^{* c}$, we have
\[
	|x - y^k| \ge \text{dist}(y^k, Q_k^{* c})
.\] 
But $B(y^k, \text{diam}Q_k) \subset Q_k^*$: if $Q_k$ is the cube centered at $y^k$ with side length $L$, then $\text{diam}(Q_k) = L\sqrt{n} $. The cube $Q_k^*$ with side length $2 \sqrt{n} L$ is the smallest cube centered at $y^k$ containing $B(y^k, \text{diam}(Q_k))$.

Hence
\[
	\text{dist}(y^k, Q_k^{* c}) \ge 
	\text{dist}(y^k, B(y^k, \text{diam}(Q_k))^c)
	= \text{diam}(Q_k) \ge 2 |y - y^k|
.\] 

(2) and (3). Fix $x \in Q_k^{* c}, y_0 \in Q_k$. Take any $y \in Q_k$. Let
\[
	\varepsilon = |x - y_0|, \varepsilon' = |x - y|, d = |x - y^k|
.\]
We have
\[
	|x - y^k| - |y^k - y_0| \le |x - y_0| \le |x - y^k| + |y^k - y_0|
.\] 
By the inequality we just proved, $|y^k - y_0| \le \frac{1}{2} d$. Thus
\[
	\frac{1}{2} d \le \varepsilon \le \frac{3}{2} d
.\]
We also have
\[
	\frac{1}{2} d \le  \varepsilon' \le \frac{3}{2} d
.\] 
Thus
\[
	\frac{1}{3} \varepsilon \le \varepsilon' \le 3 \varepsilon
.\] 

Take $r = \frac{1}{3} \varepsilon$, $R = 3 \varepsilon$, we have
\[
	Q_k \subset B(x, R) \setminus B(x, r)
.\] 
\end{proof}

\item 
Prove that
(i)

\[
e^{-\alpha}=\frac{2}{\pi} \int_0^{\infty} \frac{\cos (\alpha s)}{1+s^2} d s, \quad \forall \alpha>0
\]

(ii) Use (i), the fact that

\[
\frac{1}{1+s^2}=\int_0^{\infty} e^{-\left(1+s^2\right) t} d t
\]

and

\[
\int_0^{\infty} e^{-a^2 s^2} \cos (2 b s) d s=\frac{\sqrt{\pi}}{2 a} e^{-b^2 / a^2} \text { (Can you verify this?), }
\]

combined with Fubini's theorem, to prove that

\[
e^{-\alpha}=\frac{1}{\sqrt{\pi}} \int_0^{\infty} \frac{e^{-t}}{\sqrt{t}} e^{-\alpha^2 /(4 t)} d t
\]

\begin{proof}
	First, we note
	\[
		\int_{-\infty}^{\infty} \frac{\cos (\alpha s)}{1+s^2} d s = \int_{-\infty}^{\infty} \frac{e^{i \alpha s}}{1+s^2} d s
	.\] 
	Now, to evaluate the integral in RHS, consider it as a path integral in complex plane. We see that the integrand $\to 0$ as $\operatorname{Re} s \to \pm \infty $ and $\operatorname{Im} s \to - \infty $. Thus
	\begin{align*}
		& \int_{-\infty}^{\infty} \frac{e^{i \alpha s}}{1+s^2} d s =: \int _{- \infty } ^\infty F(s) ds \\
		&= \lim_{R \to \infty } \int _{-R}^R F(s) ds + \int_{0}^\pi F(R e^{- i \theta}) R d \theta \\
		&= - 2 \pi i \operatorname{Res} \left( F, - i \right)  \\
		&= - 2 \pi i \cdot (-\frac{e^{- \alpha}}{2 i}) = \pi e^{- \alpha}
	.\end{align*} 
	Thus $\frac{2}{\pi} \int_0^{\infty} \frac{\cos (\alpha s)}{1+s^2} d s = \frac{2}{\pi} \left( \frac{1}{2} \pi e^{- \alpha} \right) = e^{- \alpha}$.

	Now,
	\begin{align*}
		& \int_0^{\infty} \frac{\cos (\alpha s)}{1+s^2} ds\\
		&= \int_0^{\infty} \cos (\alpha s) \int _0^\infty \exp \left( -(1+s^2) t\right) d t  d s \\
		&= \int_0^{\infty} e^{- t} \int _0^\infty e^{- s^2 t}cos( \alpha s) d s  d t \\
		&= \int _0^\infty e^{- t} \frac{\sqrt{\pi} }{2 \sqrt{t} } e^{- \alpha^2 / (4 t)} \\
	.\end{align*}
	From which the result follows.
\end{proof}

\item Let

\[
P(x)=\frac{C_n}{\left(1+|x|^2\right)^{n+1 / 2}}, \quad C_n=\frac{\Gamma\left(\frac{n+1}{2}\right)}{\pi^{\frac{n+1}{2}}}
\]

be the Poisson kernel as defined as defined in class. Prove that

\[
P(x)=\int_{\mathbb{R}^n} e^{-2 \pi|\xi|} e^{-2 \pi i x \cdot \xi} d \xi
\]


Thus $\hat{P}(\xi)=e^{-2 \pi|\xi|}$.
Hint: Used the previous problem and write

\[
\int_{\mathbb{R}^n} e^{-2 \pi|\xi|} e^{-2 \pi i x \cdot \xi} d \xi=\int_{\mathbb{R}^n}\left(\frac{1}{\sqrt{\pi}} \int_0^{\infty} \frac{e^{-t}}{\sqrt{t}} e^{-4 \pi^2|\xi|^2 / 4 t} d t\right) e^{-2 \pi i x \cdot \xi} d \xi
\]

to reduce to the formula for for $\hat{H}(\xi)$.

Remark 1 This "procedure" is called the "Bochner-subordination principle." (See Stein p. 61.)

\begin{proof}
\begin{align*}
	\int_{\mathbb{R}^n} e^{-2 \pi|\xi|} e^{-2 \pi i x \cdot \xi} d \xi
	&=\int_{\mathbb{R}^n}\left(\frac{1}{\sqrt{\pi}} \int_0^{\infty} \frac{e^{-t}}{\sqrt{t}} e^{-4 \pi^2|\xi|^2 / 4 t} d t\right) e^{-2 \pi i x \cdot \xi} d \xi \\
	&=\frac{1}{\sqrt{\pi}}  \int_0^{\infty}\frac{e^{-t}}{\sqrt{t}} \int_{\mathbb{R}^n} e^{-4 \pi^2|\xi|^2 / 4 t}e^{-2 \pi i x \cdot \xi} d \xi  d t  \\
.\end{align*}
Recall
\[
	H_{\sqrt{t} }(x) = \frac{1}{\left( 4 \pi t \right) ^{n / 2}} e^{- |x|^2 / 4 t}, \quad \hat{H} _{\sqrt{t} }(\xi) = \hat{H} (\sqrt{t}  x) = e^{- 4 \pi^2 t \xi^2}
.\] 
Thus
\begin{align*}
	&=\frac{1}{\sqrt{\pi}}  \int_0^{\infty}\frac{e^{-t}}{\sqrt{t}} \int_{\mathbb{R}^n} \hat{H} _{\sqrt{\frac{1}{4 t}} }(\xi) e^{-2 \pi i x \cdot \xi} d \xi  d t  \\
	&=\frac{1}{\sqrt{\pi}}  \int_0^{\infty}\frac{e^{-t}}{\sqrt{t}} H _{\sqrt{\frac{1}{4 t}} }(t) d t  \\
	&=\frac{1}{\sqrt{\pi}}  \int_0^{\infty}\frac{e^{-t}}{\sqrt{t}} \left[ (4 \pi)^{- \frac{n}{2}} (4 t)^{n / 2} e^{- |x|^2 t} \right]  d t  \\
	&= \pi^{- (n + 1) / 2} \int _0^\infty  t^{(n - 1) / 2} e^{- (|x|^2 + 1) t} d t\\
	&= \pi^{- (n + 1) / 2} \int _0^\infty  (|x|^2 + 1)^{- (n + 1) / 2} s^{(n - 1) / 2} e^{- s} d s\\
	&= \pi^{- (n + 1) / 2} (|x|^2 + 1)^{- (n + 1) / 2} \Gamma((n+ 1) / 2)
.\end{align*}
\end{proof}

\item For $f \in L^2\left(\mathbb{R}^n\right) \cap L^1\left(\mathbb{R}^n\right)$ define

\[
u_f(x, y)=\int_{\mathbb{R}^n} \hat{f}(\xi) e^{-2 \pi y|\xi|} e^{-2 \pi i x \cdot \xi} d \xi
\]

and

\[
v_f(x, t)=\int_{\mathbb{R}^n} \hat{f}(\xi) e^{-4 \pi^2 t|\xi|^2} e^{-2 \pi i x \cdot \xi} d \xi .
\]


Prove that $u_f(x, y)=P_y * f(x)$ and $v_f(x, t)=H_t * f(x)$ and hence by previous problem $u_f$ is harmonic in $\mathbb{R}_{+}^{n+1}$ and $v_f$ satisfies the heat equation in $\mathbb{R}_{+}^{n+1}$, both with boundary values $f$. Notice, however, that these properties can be verified very easily with the above formulas by (easily justifying) differentiating under the integral sign.

\begin{proof}
	\begin{align}
\begin{aligned}
& P_y * f(x)=\int_{\mathbb{R}^n} P_y(x-z) f(z) d z=\int_{\mathbb{R}^n}\left(\int_{\mathbb{R}^n} e^{-2 \pi y|\xi|} e^{2 \pi i \xi \cdot(x-z)} d \xi\right) f(z) d z \\
& =\int_{\mathbb{R}^n} e^{-2 \pi y|\xi|} e^{2 \pi i \xi \cdot x}\left(\int_{\mathbb{R}^n} f(z) e^{-2 \pi i \xi \cdot z} d z\right) d \xi \\
& =\int_{\mathbb{R}^n} e^{-2 \pi y|\xi|} e^{2 \pi i \xi \cdot x} \hat{f}(\xi) d \xi \\
& =u_f(x, y)
\end{aligned}
\end{align}

\begin{align}
\begin{aligned}
& H_t * f(x)=\int_{\mathbb{R}^n} H_t(x-z) f(z) d z=\int_{\mathbb{R}^n}\left(\int_{\mathbb{R}^n} e^{-4 \pi^2 t|\xi|^2} e^{2 \pi i \xi \cdot(x-z)} d \xi\right) f(z) d z \\
& =\int_{\mathbb{R}^n} e^{-4 \pi^2 t|\xi|^2} e^{2 \pi i \xi \xi \cdot x} \hat{f}(\xi) d \xi \\
& =v_f(x, t)
\end{aligned}
\end{align}
\end{proof}

\item 
	Let $f \in C_0^1\left(\mathbb{R}^n\right)$. Prove
1 .

\[
\int_{\mathbb{R}^n}|\nabla f(x)|^2 d x=4 \pi^2 \int_{\mathbb{R}^n}|\xi|^2|\hat{f}(\xi)|^2 d \xi
\]

2. For all $y \in \mathbb{R}^n$ show that

\[
\int_{\mathbb{R}^n}|f(x+y)+f(x-y)-2 f(x)|^2 d x=4 \int_{\mathbb{R}^n}|\hat{f}(\xi)|^2|1-\cos (2 \pi \xi \cdot y)|^2 d \xi
\]


Remark 2 The first is the Dirichlet form of the Laplacian and the second is related to Dirichlet forms of some non-local (pseudo differential) operator-will see them a little later.

\begin{proof}
	This is proved using Plancherel's Theorem Fourier transform. 
	\begin{align*}
		\int _{\mathbb{R}^n} |\nabla f(x)|^2 dx
		&= \int _{\mathbb{R}^n} \left| \hat{\nabla  f} (\xi) \right| ^2 d \xi \\
		&= \int r_{\mathbb{R}^n} \left| \sum - 2 \pi i \xi_k \hat{f} (\xi) \right| ^2 d \xi \\
		&= \int _{\mathbb{R}^n} (2 \pi)^2 |\xi|^2 |\hat{f} (\xi)|^2 d \xi \\
		&= 4 \pi^2 \int_{\mathbb{R}^n}|\xi|^2|\hat{f}(\xi)|^2 d \xi \\
	.\end{align*}

	For the second part, we have
	\begin{align*}
		& \int_{\mathbb{R}^n}|f(x+y)+f(x-y)-2 f(x)|^2 d x \\
		&= \int_{\mathbb{R}^n} |\hat{f}(\xi) e^{-2 \pi i \xi \cdot y} + \hat{f}(\xi) e^{2 \pi i \xi \cdot y} - 2 \hat{f}(\xi)|^2 d \xi \\
		&= \int_{\mathbb{R}^n} |\hat{f}(\xi)|^2 |e^{-2 \pi i \xi \cdot y} + e^{2 \pi i \xi \cdot y} - 2|^2 d \xi \\
		&= \int_{\mathbb{R}^n}|\hat{f}(\xi)|^2|2- 2\cos (2 \pi \xi \cdot y)|^2 d \xi \\
		\end{align*}
\end{proof}


\item Let $f \in L^2\left(\mathbb{R}^n\right)$ and set $u_f(x, y)=P_y * f(x)$ and $v_f(x, t)=H_t * f(x)$. Set

\[
\nabla_x u_f=\left(\frac{\partial u_f}{\partial x_1}, \frac{\partial u_f}{\partial x_2}, \ldots, \frac{\partial u_f}{\partial x_n}\right)
\]


Define the Littlewood-Paley functions (will return to them later)

\[
\begin{aligned}
g_v(f)(x) & =\left(\int_0^{\infty} y\left|\frac{\partial u_f}{\partial y}(x, y)\right|^2 d y\right)^{1 / 2} \\
g_h(f)(x) & =\left(\int_0^{\infty} y\left|\nabla_x u_f(x, y)\right|^2 d y\right)^{1 / 2} \\
g(f)(x) & =\left(\int_0^{\infty} y\left|\nabla u_f(x, y)\right|^2 d y\right)^{1 / 2} \\
G_v(f)(x) & =\left(\int_0^{\infty} t\left|\frac{\partial v_f}{\partial t}(x, t)\right|^2 d t\right)^{1 / 2} \\
G_h(f)(x) & =\left(\int_0^{\infty}\left|\nabla_x v_f(x, t)\right|^2 d t\right)^{1 / 2}
\end{aligned}
\]


Prove that

\[
\left\|g_v(f)\right\|_2=\left\|g_h(f)\right\|_2=\left\|G_v(f)\right\|_2=\frac{1}{2}\|f\|_2
\]

and that

\[
\left\|G_h(f)\right\|_2=\frac{1}{\sqrt{2}}\|f\|_2
\]

\begin{proof}
	We have
	\begin{align*}
		\left\|g_v(f)\right\|_2^2
		&= \int _{\mathbb{R}^n} \left( \int _0^\infty y \left| \frac{\partial u_f}{\partial y} (x, y) \right| ^2 dy \right) dx \\
		&= \int_0 ^\infty y \int _{\mathbb{R}^n} \left| \frac{\partial u_f}{\partial y} (x, y) \right| ^2 dx dy \\
		&= \int _0^\infty y \int _{\mathbb{R}^n} \left| \widehat{ \frac{\partial u_f}{\partial y} }(\xi)\right| ^2 d \xi dy \\
	\end{align*}
	Now, note that
	\[
		\widehat{ \frac{\partial u_f}{\partial y} }(\xi) = \int \frac{\partial u_f}{\partial y} (x, y) e^{2 \pi i x \cdot \xi} dx = \frac{\partial}{\partial y} \left( \hat{P_y}(\xi) \hat{f}(\xi)
		\right) 
	.\]
	\[
		\hat{P_y} (\xi) = e^{-2 \pi y |\xi|}
	\]
	\[
			\int _0^\infty y \left| \frac{\partial}{\partial y} \left( \hat{P_y}(\xi) \right) \right| ^2 d y = \int _0^\infty y (4 \pi^2 |\xi|^2) e ^{-4 \pi y |\xi|} dy = \frac{1}{4}
	\]
	Thus
	\begin{align*}
		\left\|g_v(f)\right\|_2^2
		&= \int _0^\infty y \int _{\mathbb{R}^n} \left| \frac{\partial}{\partial y} \left( \hat{P_y}(\xi) \right) \right| ^2 |\hat{f}(\xi)|^2 d \xi dy \\
		&= \int_ {\mathbb{R}^n} \int _0^\infty y \left| \frac{\partial}{\partial y} \left( \hat{P_y}(\xi) \right) \right| ^2 d y |\hat{f}(\xi)|^2 d \xi \\
		&= \int _{\mathbb{R}^n} \frac{1}{4} |\hat{f}(\xi)|^2 d \xi \\
		&= \frac{1}{4} \|f\|_2^2
	\end{align*}
	
	Next, we control $\left\|g_h(f)\right\|_2$.

	\begin{align*}
		\left\|g_h(f)\right\|_2^2
		&= \int _{\mathbb{R}^n} \left( \int _0^\infty y \left| \nabla_x u_f(x, y) \right| ^2 dy \right) dx \\
		&= \int_0 ^\infty y \int _{\mathbb{R}^n} \left| \nabla_x u_f(x, y) \right| ^2 dx dy \\
		&= \int _0^\infty y \int _{\mathbb{R}^n} \left| \nabla_x \widehat{u_f}(\xi, y) \right| ^2 d \xi dy \\
		&= \int _0^\infty y \int _{\mathbb{R}^n} \left| \widehat{\nabla_x  P_y(\xi)} \hat{f}(\xi) \right| ^2 d \xi dy \\
		&= \int _0^\infty y \int _{\mathbb{R}^n} 4 \pi^2 |\xi|^2 \left| P_y(\xi) \hat{f}(\xi)  \right| ^2 d \xi dy \\
	\end{align*}

	where

	\[
		\int_0^\infty y \widehat{P}_y(\xi) ^2 dy = \int_0^\infty y e^{-4 \pi y |\xi|} dy = \frac{1}{16 \pi^2 |\xi|^2}
	.\]

	Thus
	\begin{align*}
		\left\|g_h(f)\right\|_2^2
		&= \frac{1}{4} \|\hat{f}\|_2^2 \\
		&= \frac{1}{4} \|f\|_2^2
	\end{align*}

	The proof for $\left\|G_v(f)\right\|_2$ is similar to the proof for $\left\|g_v(f)\right\|_2$. We only need to replace $P_y$ with $H_t$ and $\nabla_x$ with $\frac{\partial}{\partial t}$. Use the following:

	\begin{align*}
		&\int_{0}^{\infty} t \left| \frac{\partial}{\partial t} \left( \hat{H_t}(\xi) \right) \right| ^2 d t \\
		&= \int_{0}^{\infty} t \left| \frac{\partial}{\partial t} \left( e^{-4 \pi^2 t |\xi|^2} \right) \right| ^2 d t \\
		&= \int_{0}^{\infty} t \left| - 4 \pi^2 |\xi|^2 e^{-4 \pi^2 t |\xi|^2} \right| ^2 d t \\
		&= \int_{0}^{\infty} 16 \pi^4 |\xi|^4 t e^{-8 \pi^2 t |\xi|^2} d t \\
		&= \frac{1}{4}
	\end{align*}

	Finally, for $\left\|G_h(f)\right\|_2$, we have	

	\[
		 \widehat{\nabla_x v_f}(\xi, t) 
		 = \widehat{\nabla_x H_t}(\xi) \hat{f}(\xi)
	\]


	\begin{align*}
		\| G_h(f) \|_2^2
		&= \int _{\mathbb{R}^n} \left( \int _0^\infty t \left| \nabla_x v_f(x, t) \right| ^2 dt \right) dx \\
		&= \int_0 ^\infty t \int _{\mathbb{R}^n} \left| \nabla_x v_f(x, t) \right| ^2 dx dt \\
		&= \int _0^\infty t \int _{\mathbb{R}^n} \left|  \widehat{\nabla_x H_t}(\xi) \hat{f}(\xi) \right| ^2 d \xi dt \\
	\end{align*}
	
	where
	\begin{align*}
		| \widehat{\nabla_x H_t}(\xi) |^2
		&= \left( \nabla_x H_t * \nabla_x H_t(\xi) \right)^\wedge \\
		&= (\Delta_x \left( H_t *  H_t \right))^\wedge (\xi) \\
		&= (\Delta_x H_{2 t} )^\wedge (\xi) \\
		&= -4 \pi^2 |\xi|^2 \hat{H}_{2 t}(\xi) 
	.\end{align*}

	Hence
	\begin{align*}
		\| G_h(f) \|_2^2
		&= \int _0^\infty t \int _{\mathbb{R}^n} 4 \pi^2 |\xi|^2 \hat{H}_{2 t}(\xi) \left| \hat{f}(\xi) \right| ^2 d \xi dt \\
		&= \int_ {\mathbb{R}^n} 4 \pi^2 |\xi|^2 \left[
		\int _0^\infty t \hat{H}_{2 t}(\xi) d t \right] \left| \hat{f}(\xi) \right| ^2 d \xi \\
	\end{align*}
	We have
	\begin{align*}
		\int _0^\infty t \hat{H}_{2 t}(\xi) d t
		&= \int _0^\infty t e^{-4 \pi^2 (4 t^2) |\xi|^2} d t \\
		&= \frac{1}{32 \pi^2 |\xi|^2}
	\end{align*}
	Therefroe

	\[
		\| G_h(f) \|_2^2 = \frac{1}{8} \| f \|_2^2			
	.\]

		

\end{proof}

\item Let $f, h \in L^2\left(\mathbb{R}^n\right)$. Prove that

\[
\int_0^{\infty} \int_{\mathbb{R}^n} y \frac{\partial u_f}{\partial y}(x, y) \frac{\overline{\partial u_h}}{\partial y}(x, y) d x d y=\frac{1}{4} \int_{\mathbb{R}^n} f(x) \overline{h(x)} d x
\]

and

\[
\int_0^{\infty} \int_{\mathbb{R}^n} y \nabla_x u_f(x, y) \cdot \overline{\nabla_x u_h}(x, y) d x d y=\frac{1}{4} \int_{\mathbb{R}^n} f(x) \overline{h(x)} d x
\]


Remark 3 Later in the course (see Chapter 10) we will derive $L^P$-versions of the above formulas . For now, let us observe that

\[
\left|\nabla u_f\right|^2=\frac{1}{2} \Delta u_f^2
\]

and hence it follows from this exercise that

\[
\int_0^{\infty} \int_{\mathbb{R}^n} y \Delta u_f^2(x, y) d x d y=\int_{\mathbb{R}^n} f(x)^2 d x
\]

for all $f \in L^2\left(\mathbb{R}^n\right)$.

\begin{proof}
	We have
	\begin{align*}
		& \int_0^{\infty} \int_{\mathbb{R}^n} y \frac{\partial u_f}{\partial y}(x, y) \frac{\overline{\partial u_h}}{\partial y}(x, y) d x d y \\
		&= \int_0^{\infty} y \int_{\mathbb{R}^n} \widehat{ \frac{\partial u_f}{\partial y}}(\xi, y) \overline{\widehat{ \frac{\partial u_h}{\partial y}}}(\xi, y) d \xi d y \\	
		&= \int_0^{\infty} y \int_{\mathbb{R}^n} \left[ \frac{\partial}{\partial y} \hat{P}_y(\xi) \hat{f}(\xi) \right] \left[ \frac{\partial}{\partial y} \overline{\hat{P}_y(\xi) \hat{h}(\xi)} \right] d \xi d y \\
		&= \int_{\mathbb{R}^n} \left[ \int_0^{\infty} y \left| \frac{\partial}{\partial y} \hat{P}_y(\xi) \right|^2 d y \right] \hat{f}(\xi) \overline{\hat{h}(\xi)} d \xi \\	
	\end{align*}
	where
	\[
		\int_0^{\infty} y \left| \frac{\partial}{\partial y} \hat{P}_y(\xi) \right|^2 d \xi = \int_0^{\infty} y (4 \pi^2 |\xi|^2) e^{-4 \pi y |\xi|} dy = \frac{1}{4}
	.\]	
	Thus
	\begin{align*}
		\int_0^{\infty} \int_{\mathbb{R}^n} y \frac{\partial u_f}{\partial y}(x, y) \frac{\overline{\partial u_h}}{\partial y}(x, y) d x d y 
		&= \frac{1}{4} \int_{\mathbb{R}^n} \hat{f}(\xi) \overline{\hat{h}(\xi)} d \xi \\
		&= \frac{1}{4} \int_{\mathbb{R}^n} f(x) \overline{h(x)} d x
	\end{align*}

	For the second part, we have

	\begin{align*}
		& \int_0^{\infty} \int_{\mathbb{R}^n} y \nabla_x u_f(x, y) \cdot \overline{\nabla_x u_h}(x, y) d x d y \\
		&= \int_0^{\infty} y \int_{\mathbb{R}^n} \widehat{\nabla_x u_f}(\xi, y) \cdot \overline{\widehat{\nabla_x u_h}}(\xi, y) d \xi d y \\
		&= \int_0^{\infty} y \int_{\mathbb{R}^n} \left[ \widehat{\nabla_x P_y}(\xi) \hat{f}(\xi) \right] \cdot \left[ \overline{\widehat{\nabla_x P_y}(\xi) \hat{h}(\xi)} \right] d \xi d y \\
		&= \int_{\mathbb{R}^n} \left[ \int_0^{\infty} y \widehat{\Delta_x {P}_{2y}} (\xi) d y \right] \hat{f}(\xi) \overline{\hat{h}(\xi)} d \xi \\
		&= \int_{\mathbb{R}^n} (4 \pi^2 |\xi|^2) \left[ \int_0^{\infty} y \hat{P}_{2y} (\xi) d y \right] \hat{f}(\xi) \overline{\hat{h}(\xi)} d \xi \\
		&= \int_{\mathbb{R}^n} \frac{1}{4} \hat{f}(\xi) \overline{\hat{h}(\xi)} d \xi \\
		&= \frac{1}{4} \int_{\mathbb{R}^n} f(x) \overline{h(x)} d x
	\end{align*}

	where
	\begin{align*}
		\int_0^{\infty} y \hat{P}_{2y} (\xi) d y
		&= \int_0^{\infty} y e^{-2 \pi (2 y) |\xi|} dy \\
		&= \frac{1}{16 \pi^2 |\xi|^2}
	\end{align*}
\end{proof}

\item Let $\varphi$ be a $C_0^{\infty}\left(\mathbb{R}^n\right)$ radial function. As shown in class, its Fourier transform is also radial. Suppose

\[
\int_0^{\infty}\left|\hat{\varphi}^{\prime}(y)\right|^2 y d y=A<\infty
\]


Define

\[
U(x, y)=\left(\varphi_y * f\right)(x)
\]

and consider the generalized Littlewood-Paley function

\[
g(f)(x)=\left(\int_0^{\infty} y\left|\frac{\partial U}{\partial y}(x, y)\right|^2 d y\right)^{1 / 2}
\]


Prove that

\[
\|g(f)\|_2=\sqrt{A}\|f\|_2
\]

\begin{proof}
	We have
	\begin{align*}
		\|g(f)\|_2^2
		&= \int _{\mathbb{R}^n} \left( \int _0^\infty y \left| \frac{\partial U}{\partial y} (x, y) \right| ^2 dy \right) dx \\
		&= \int_0 ^\infty y \int _{\mathbb{R}^n} \left| \frac{\partial U}{\partial y} (x, y) \right| ^2 dx dy \\
		&= \int _0^\infty y \int _{\mathbb{R}^n} \left| \frac{\partial}{\partial y} \left( \varphi_y * f \right) (x, y) \right| ^2 dx dy \\
		&= \int _0^\infty y \int _{\mathbb{R}^n} \left| \left( \frac{\partial}{\partial y} \varphi_y * f \right) (x, y) \right| ^2 dx dy \\
		&= \int _0^\infty y \int _{\mathbb{R}^n} \left| \widehat{\varphi_y'}(\xi) \hat{f}(\xi) \right| ^2 d \xi dy \\
		&= \int _{\mathbb{R}^n} \int _0^\infty y \left| \widehat{\varphi_y'}(\xi) \right| ^2 d y \left| \hat{f}(\xi) \right| ^2 d \xi \\	
		&= A \| \hat{f} \|_2^2 \\
		&= A \| f \|_2^2
	\end{align*}
\end{proof}

\item Let $f$ be a function on the real line $\mathbb{R}$ such that both $f$ and $x f$ are in $L^2(\mathbb{R})$. Prove that $f \in L^1(\mathbb{R})$ and that

\[
\left(\int_{\mathbb{R}}|f(x)| d x\right)^2 \leq 8\left(\int_{\mathbb{R}}|f(x)|^2 d x\right)^{1 / 2}\left(\int_{\mathbb{R}}|x|^2|f(x)|^2 d x\right)^{1 / 2}
\]

\begin{proof}
	We can show that $f \in L^1(\mathbb{R})$ by the following:
	\begin{align*}
		\int_{\mathbb{R}} |f(x)| dx
		&= \int_{\mathbb{R}} (f(x) (1 + |x|^2)^{1 / 2}) \frac{1}{(1 + |x|^2)^{1 / 2}} dx \\
		&\le \left( \int_{\mathbb{R}} |f(x)|^2 (1 + |x|^2) dx \right) ^{1 / 2} \left( \int_{\mathbb{R}} \frac{1}{1 + |x|^2} dx \right) ^{1 / 2}
	\end{align*}
	The left hand side equals $(\| f \|_2^2 + \| x f \|_2^2)^{1 / 2}$, and the right hand side is finite. Thus $f \in L^1(\mathbb{R})$.

	Now, to prove the inequality, we use the strategy for proving ``Peter-Paul inequality''.	

	\begin{align*}
		\int_{\mathbb{R}} |f(x)| dx 
	&=  \int_{|x| > \lambda} |f(x)| dx + \int_{|x| \le \lambda} |f(x)| dx \\	
	&=  \int_{|x| > \lambda} |xf| |x|^{-1} dx + \int_{|x| \le \lambda} |f(x)| \cdot 1 dx \\
	&\le \left( \int_{|x| > \lambda} |x f|^2 dx \right) ^{1 / 2} \left( \int_{|x| > \lambda} |x|^{-2} dx \right) ^{1 / 2} + \lambda \| f \|_2 \sqrt{2 \lambda} \\
	&\le \| x f\| _2 \sqrt{2 / \lambda} + \| f \|_2 \sqrt{2 \lambda} 
	\end{align*}	
	Choose $\lambda = \| x f \|_2 / \| f \|_2$ to get the desired inequality.
\end{proof}


\end{enumerate}
